\par\Needspace{6\baselineskip}\medskip
\hypertarget{native-142-article103}{}\pdfbookmark[2]{103.}{article-103}
\begin{center}103.\end{center}
Wenn $z$ einen größeren Wert außerhalb der Grenzen von Tafel III annimmt, lassen sich die Gleichungen 13 und $13^*$ stets zuverlässig und bequem in unveränderter Form durch Probieren lösen, und zwar aus ähnlichen Gründen, wie wir sie in Art. 94 für die Ellipse dargelegt haben. In einem solchen Fall kann man \SourceLeaf{143}voraussetzen, dass die Bahnelemente wenigstens ungefähr bekannt sind. Ein Näherungswert von $n$ ergibt sich dann sofort aus der Formel $\operatorname{tang}2n=\dfrac{\sin f\mathpunct{.}\BareRoot{rr'}}{\alpha\BareRoot{(ee-1)}}$, die unmittelbar aus Gleichung 6 in Art. 99 folgt. Aus $n$ erhält man $z$ mit der Formel $z=\dfrac{1-\cos2n}{2\cos2n}=\dfrac{\sin n^2}{\cos2n}$. Vom Näherungswert von $z$ ausgehend findet man durch wenige Versuche den Wert, der Gleichung 13 oder $13^*$ genau erfüllt. Diese Gleichungen lassen sich auch in der Form schreiben
\[
\begin{aligned}
m={}&\left(l-\frac{\sin n^2}{\cos2n}\right)^{\frac12}\\[.6em]
&+2\left(l-\frac{\sin n^2}{\cos2n}\right)^{\frac32}\frac{\dfrac{\operatorname{tang}2n}{\cos2n}-\log\operatorname{hyp}\operatorname{tang}(45^0+n)}{\operatorname{tang}2n^3}
\end{aligned}
\]
\[
\begin{aligned}
M={}&-\left(L+\frac{\sin n^2}{\cos2n}\right)^{\frac12}\\[.6em]
&+2\left(L+\frac{\sin n^2}{\cos2n}\right)^{\frac32}\frac{\dfrac{\operatorname{tang}2n}{\cos2n}-\log\operatorname{hyp}\operatorname{tang}(45^0+n)}{\operatorname{tang}2n^3}
\end{aligned}
\]
\noindent und so lässt sich unter Weglassung von $z$ der wahre Wert von $n$ unmittelbar ermitteln.
